% ======================= 第 七 章 =======================
\chapter{月兔之乱}
\chaptersubtitle{ГЛАВА СЕДЬМАЯ}{МЯТЕЖ ЛУННЫХ ЗАЙЦЕВ}
\epigraphbox{“兽穷则啮，鸟穷则啄。”}{《史记·淮阴侯列传》}

\kaishou{这场乱事在西方的文献里有另一个名字，叫月面战争。本章要说明的是：交战的双方其实是同一批造物，而镇压用的东西，是这座城自己的存货。}

\noindent
第六年，一位驻地球主任被隔离。这一条记载本身很短，短到不像一件大事。

它引发的后果，是这座城两万年间唯一一次真正的动乱。动乱之后，城的秩序还在，人口、系统、编制都在，可是它已经不是一个能继续过下去的地方。

本章按起因、经过、结局、之后四段写。四段里，只有第一段有材料，其余三段全靠影像与推算。

\section{统属中断的七年}

月兔的统属关系一向简单：归一个总管调遣。前二〇七〇年起，这个总管是嫦娥；她被隔离之后，名册仍挂在她名下，只是本人不在。

这一条在工籍上的表现很具体。残卷 `ФВ-1968-14` 第 5 叶记的是隔离当年之后第七年的编制，班次照旧，定额照旧，只在总管一栏写着「缺」，缺了七年。

七年之后，班次开始不对。第 6 叶的轮值表里有三处空缺，填补的是邻班，注记写「暂代」。第 8 叶的注记变成了「无人」。

也就是说，这件事的性质在前十年里变了三次：先是没人管，再是借人管，最后是没人可借。

造物的问题是：它们没有被教过怎么在没有人管的时候继续。它们被造出来是要听指令的，指令停了，它们能做的是等。等到某一刻，等不下去了。

那一刻发生了什么，档案不记。笔者能确认的只有一件事：那一年之后，城的记录里出现了「乱」这个字。

\section{兔对兔的四年}

关于经过，材料只有一句：这一战自始即为兔对兔。

这句话里的信息量很大。\footnote{原文见档案 `ФВ-1969-02` 第 11 叶；页边另有铅笔一行，字迹与正文不同，内容为「查过」。}月之都没有军队，唯一可称为武装的力量是月之使者，而月之使者本身也是月兔。所以当骚动起来的时候，城能派出去的镇压力量，与骚动的那一方是同一个物种。

数量对比是外人难以想象的。月之使者据说不足十只，而对方有一百多万。

前三年是拉锯。城里的月之民不参加战斗，他们不能参加：他们的身体为无秽而设计，没有再生能力，受伤即永久受损。他们的办法是关环门、停气路、断水，把骚动限制在若干环内。

这个办法在第一年有效，第三年开始失效。原因是月兔熟悉这些门。它们修过门、清过管道、送过东西，城里的每一道缝它们都走过。

第四年以后，城里的记录换了一种写法：「某环不报」。不报，就是丢了。

\section{净秽山的开启}

月之民最后做了一件事：开启净秽山。

这件事的操作细节见第二章所论的那份总目与本章所引的 `ФВ-1969-02`。此处只说结果：山里的存货被排出来，再经一次处理，还原为活性污秽；存量在短时期内放尽，城里的月兔因此全灭。

执行这一步的是月夜见本人。他进山操作，没有出来。处理产生的污秽把他灌在里面，月之都的领袖死在这一步上。那一年，他大约一千九百岁。

同一年死掉的月之民，共两万二千。这个数字记在事故一栏，没有分项。

污秽对月之民属于外伤：短时接触致死或致伤，事后不留痕迹，不会转为长期状态。本次事件没有一例属于后者。

这是这座城历史上唯一一次对自己使用污秽。

西方的文献把这件事叫作月面战争。\footnote{该名在英文文献中作 \textit{Lunar War}，中文材料一律不用；本书从西名，仅为便于与那批文献对读。}这个叫法在一点上是准确的：它确实是一场战争，双方都有战术、都有伤亡、都有战线。不准确的地方在于「双方」二字：那两方是同一批造物，一批被教过服从，一批被教过执法。

\begin{archive}[绝密／СОВЕРШЕННО СЕКРЕТНО]{档案 № ФВ-1969-02}
事故统计，共七叶，第一叶为汇总，其余按环分列。

汇总叶记：事故之数，兔乱二万二千。

第二至第七叶分记六环，每环一叶，叶内列「乱」「伤」「亡」三栏。六叶的「亡」栏相加，与汇总叶的数目完全一致；「乱」与「伤」两栏相加则为零。

第七叶末行有墨涂，涂痕下可辨「不可使聚」四字，与 `ФВ-1968-14` 第 11 叶的涂痕同字。
\end{archive}

\section{自动运转的一千七百六十年}

月兔全灭之后，城的运转没有立刻停止。

系统仍在自动运行：供能、供水、控温照旧，抽取污秽的那一段也还在转。变的是入库的东西：没有人再把污秽加工成蓬莱药，此后入库的是原状污秽，不稳定，会扩散，需要持续压制。

也就是说，这座城在失去全部劳力之后，仍然自动运转了一千七百六十年。

这一千七百六十年里，城里的人做三件事：压制山里的污秽、维持剩下的系统、等待。等到第四批人下去之后，等待变成了准备。

撤离前的清点见第八章。此处只说一句：那一千七百六十年间，城的档案里没有一条记录提到「重建」。这个动作不存在：第二章所说的那种判断，此时已经停了两万年。

\section{本章结论}

本章说了四件事。

起因是没人管。总管一栏缺了七年，班次错了三次，最后无人可借。

经过是兔对兔。城能派的力量与骚动的一方同种，数量相差五个数量级。

结局是开山。用存货换取全灭，代价是两万二千个自己人与一位领袖。

之后是等。系统自动运转一千七百六十年，没有人提议重建。

四件事合起来，说明的是本书的总题目在这一点上如何兑现：他们把「变化」删掉了，于是没有新人口，于是造了会生育的劳工，于是劳工的数量失控，于是只能用存货处理。每一步都是前一步的直接后果，没有一步是意外。

本章的自我拆台：这是一场一方的记录。所有材料都由月之民书写，月兔那一方一个字也没有留下，连它们为什么在第七年之后开始不等了，都没有记录。本章写「等不下去了」，是从班次表推的，不是从它们的话推的。

严格地说，本章只写了胜利者的账。而胜利者在这一战之后，已经不是原来那个胜利者了。
